\documentclass[10pt, a4paper]{article}

% --- PACKAGES FOR COMPACT LAYOUT ---
\usepackage[margin=0.5in]{geometry} % Tight margins
\usepackage{multicol}               % Two-column layout
\usepackage[compact]{titlesec}      % Reduces space around section titles
\usepackage{enumitem}               % For compact lists

% --- PACKAGES FOR CONTENT ---
\usepackage{listings}               % Code snippets with syntax highlighting
\usepackage{xcolor}                 % Custom colors
\usepackage[most]{tcolorbox}        % Colored boxes
\usepackage{booktabs}               % Professional tables
\usepackage{fontawesome5}           % Icons for flair

% --- COLOR DEFINITIONS ---
\definecolor{codegreen}{rgb}{0,0.6,0}
\definecolor{codegray}{rgb}{0.5,0.5,0.5}
\definecolor{codepurple}{rgb}{0.58,0,0.82}
\definecolor{backcolour}{rgb}{0.95,0.95,0.92}
\definecolor{tipboxcolor}{rgb}{0.9, 0.98, 1.0} % Light blue

% --- LISTINGS (CODE) STYLE FOR PYTHON ---
\lstdefinestyle{pythonstyle}{
    backgroundcolor=\color{backcolour},   
    commentstyle=\color{codegreen},
    keywordstyle=\color{magenta},
    numberstyle=\tiny\color{codegray},
    stringstyle=\color{codepurple},
    basicstyle=\ttfamily\footnotesize,
    breakatwhitespace=false,         
    breaklines=true,                 
    captionpos=b,                    
    keepspaces=true,                 
    numbers=left,                    
    numbersep=5pt,                  
    showspaces=false,                
    showstringspaces=false,
    showtabs=false,                  
    tabsize=2,
    language=Python
}
\lstset{style=pythonstyle}

% --- DOCUMENT LAYOUT CONFIGURATION ---
\pagestyle{empty} % No page numbers
\setlength{\columnsep}{1cm} % Space between columns
\titlespacing*{\section}{0pt}{*1.5}{*1.5} % Compact section spacing
\titlespacing*{\subsection}{0pt}{*1}{*1} % Compact subsection spacing
\setlist{noitemsep, topsep=2pt, leftmargin=*} % Compact list settings

% --- BEGIN DOCUMENT ---
\begin{document}
\begin{center}
    {\Huge\bfseries Python Quick Reference} \\
    {\large A Cheat Sheet for Core Concepts}
\end{center}

\begin{multicols}{2} % Start two-column layout

% --- SECTION: BASIC SYNTAX & DATA TYPES ---
\section*{Basics & Data Types}
Variables are dynamically typed. Indentation (4 spaces) is mandatory. Use `#` for comments.
\begin{lstlisting}
# Variable assignment
name = "Alice"  # String
age = 30        # Integer
height = 5.5    # Float
is_student = True # Boolean

# Print output with an f-string
print(f"{name} is {age} years old.")
\end{lstlisting}

% --- SECTION: DATA STRUCTURES ---
\section*{Data Structures}
\subsection*{Lists \faList}
Ordered, mutable sequences.
\begin{lstlisting}
# List creation and indexing
nums = [10, 20, 30, 40, 50]
print(nums[0])      # -> 10
print(nums[-1])     # -> 50
print(nums[1:4])    # -> [20, 30, 40]

# List methods
nums.append(60)     # Add to end
nums.pop(2)         # Remove at index 2
\end{lstlisting}

\subsection*{Dictionaries \faBook}
Unordered key-value pairs.
\begin{lstlisting}
# Dictionary creation
person = {"name": "Bob", "age": 25}

# Accessing values
print(person["name"]) # -> Bob
print(person.get("city", "N/A"))

# Modifying
person["age"] = 26
person["city"] = "New York"
\end{lstlisting}

% --- SECTION: OPERATORS ---
\section*{Common Operators}
\begin{tabular}{@{}ll@{}}
\toprule
\textbf{Category} & \textbf{Operators} \\
\midrule
Arithmetic & \texttt{+, -, *, /, \%, **, //} \\
Comparison & \texttt{==, !=, <, >, <=, >=} \\
Logical & \texttt{and, or, not} \\
Assignment & \texttt{=, +=, -=, *=, /=} \\
Membership & \texttt{in, not in} \\
\bottomrule
\end{tabular}
\vfill\null % Distribute space within the column
\columnbreak % Manually break to the next column

% --- SECTION: CONTROL FLOW ---
\section*{Control Flow}
\subsection*{Conditional Statements}
\begin{lstlisting}
score = 85
if score >= 90:
    print("Grade: A")
elif score >= 80:
    print("Grade: B")
else:
    print("Grade: C")
\end{lstlisting}

\subsection*{Loops}
\begin{lstlisting}
# For loop (iterating over a list)
fruits = ["apple", "banana", "cherry"]
for fruit in fruits:
    print(fruit)

# While loop
count = 0
while count < 3:
    print(f"Count: {count}")
    count += 1
\end{lstlisting}

% --- SECTION: FUNCTIONS ---
\section*{Functions}
\subsection*{Defining a Function}
\begin{lstlisting}
def greet(name, msg="Hello"):
    """Greets a person."""
    return f"{msg}, {name}!"

# Calling the function
print(greet("Eve")) # -> Hello, Eve!
\end{lstlisting}

\subsection*{Lambda Functions}
\begin{lstlisting}
# A lambda that adds 10 to a number
add_ten = lambda x: x + 10
print(add_ten(5)) # -> 15
\end{lstlisting}

% --- SECTION: FILE I/O ---
\section*{File I/O \faFileAlt}
Use `with` for safe file handling.
\begin{lstlisting}
# Writing to a file
with open("data.txt", "w") as f:
    f.write("Hello, World!\n")

# Reading from a file
with open("data.txt", "r") as f:
    content = f.read()
\end{lstlisting}

% --- SECTION: IMPORTANT TIPS BOX ---
\begin{tcolorbox}[
    colback=tipboxcolor, colframe=blue!50!black,
    title=\bfseries\faLightbulb~Pythonic Tips,
    fonttitle=\bfseries, arc=2mm, boxrule=1pt, top=2pt, bottom=2pt
]
\begin{itemize}
    \item \textbf{List Comprehensions:} Concise list creation. \\
    \texttt{squares = [x**2 for x in range(5)]}
    \item \textbf{f-Strings:} Modern string formatting. \\
    \texttt{name="AI"; print(f"Hello, {name}!")}
    \item \textbf{Virtual Environments:} Isolate project dependencies. \\
    \texttt{python -m venv myenv}
    \item \textbf{Unpacking:} Assign iterable items to variables. \\
    \texttt{a, b, c = [1, 2, 3]}
\end{itemize}
\end{tcolorbox}

\end{multicols} % End two-column layout
\end{document}